\section{Further structure of the exact support shadow}
\label{app:exact-shadow-details}

This appendix records exact-support structure that is useful for auditing and extending the support shadow but is not required for the main narrative.

\subsection{Support transport}

For finite discrete spaces define
\[
\supp\mu:=\{x:\mu(x)>0\}.
\]

\begin{proposition}[Finite support shadow]
\label{prop:app-support-shadow-functor}
If $f:(X,\mu)\to(Y,\nu)$ is a morphism in the finite discrete subcategory of $\int\Prob$, then
\[
\supp\nu=f(\supp\mu).
\]
Consequently support defines a functor from finite state-decorated spaces to finite sets by restricting each deterministic map to the support of its source state.
\end{proposition}

\begin{proof}
For $y\in Y$,
\[
\nu(y)
=(f_*\mu)(y)
=\sum_{x:f(x)=y}\mu(x),
\]
which is positive exactly when some $x\in\supp\mu$ satisfies $f(x)=y$.
\end{proof}

\subsection{Measure-class invariance of exact realization}

Let
\[
\Rcal^{\mathrm{a.s.}}_\Phi(\mu)
:=
\{K\subseteq J:\exists h:Z_K\to Y\text{ measurable with }\Phi=hQ_K\ \mu\text{-a.s.}\}.
\]

\begin{proposition}[Exact localization factors through measure class]
\label{thm:app-measure-class-invariance}
If $\mu\sim\nu$, then
\[
\Rcal^{\mathrm{a.s.}}_\Phi(\mu)
=
\Rcal^{\mathrm{a.s.}}_\Phi(\nu).
\]
\end{proposition}

\begin{proof}
For fixed $K$ and $h$, let
\[
B_{K,h}:=\{x:\Phi(x)\neq h(Q_K(x))\}.
\]
The decoder realizes $\Phi$ exactly under $\mu$ iff $B_{K,h}$ is $\mu$-null, which under $\mu\sim\nu$ is equivalent to being $\nu$-null. Existence of a realizing decoder is therefore the same under both laws.
\end{proof}

Together with Proposition~\ref{prop:app-measure-class-pushforward}, this gives a process-level coarse-graining: equivalent source states remain equivalent at every deterministic cut, and exact almost-sure localization cannot distinguish their different positive weights.

\subsection{Support restriction and truth-valued localization}

For finite $S\subseteq X$, write
\[
\Rcal_\Phi(S)
:=
\{K\subseteq J:\exists h:Q_K(S)\to Y\text{ with }\Phi|_S=hQ_K|_S\}.
\]

\begin{proposition}[Support restriction]
\label{thm:app-support-restriction-shadow}
If $S\subseteq T$, then
\[
\Rcal_\Phi(T)\subseteq\Rcal_\Phi(S).
\]
\end{proposition}

\begin{proof}
Any decoder that works on $T$ still works after restriction to $S$.
\end{proof}

Write
\[
\Supp(X):=(\mathcal P(X),\subseteq),
\qquad
\Two=\{0<1\}.
\]
Exact support localization can be represented as
\begin{equation}
\Loc_{\Phi,Q}:
\Supp(X)^{\mathrm{op}}\times\Bool_J
\longrightarrow\Two,
\label{eq:app-support-shadow-profunctor}
\end{equation}
with $\Loc_{\Phi,Q}(S,K)=1$ exactly when $K\in\Rcal_\Phi(S)$. Currying gives
\begin{equation}
\mathscr R_{\Phi,Q}:
\Supp(X)^{\mathrm{op}}
\longrightarrow
\Up(\Bool_J),
\qquad
S\longmapsto\Rcal_\Phi(S).
\label{eq:app-support-shadow-localization-functor}
\end{equation}
This truth-valued object remembers exact factorization on support and forgets all probability weights.

\subsection{Discernibility clutters and blockers}

For a map $f:S\to T$, write
\[
\Eq(f)
:=
\{(x,x')\in S^2:f(x)=f(x')\}.
\]
For a phenomenon-distinguished pair $x,x'\in S$, define
\[
E(x,x')
:=
\{j\in J:Q_j(x)\neq Q_j(x')\},
\]
and the distinction hypergraph
\[
\Hcal_\Phi(S)
:=
\{E(x,x'):x,x'\in S,\ \Phi(x)\neq\Phi(x')\}.
\]
For a hypergraph $\Hcal\subseteq2^J$, let
\[
\Hit(\Hcal)
:=
\{K\subseteq J:K\cap E\neq\varnothing\text{ for every }E\in\Hcal\}.
\]
The main text records the equivalence
\[
K\in\Rcal_\Phi(S)
\Longleftrightarrow
\Eq(Q_K|_S)\subseteq\Eq(\Phi|_S)
\Longleftrightarrow
K\text{ hits every }E\in\Hcal_\Phi(S).
\]

Deleting hyperedges that strictly contain another gives the clutter
\[
\Ccal_\Phi(S):=\Min_{\subseteq}\Hcal_\Phi(S).
\]
If
\[
b(\Ccal):=\Min_{\subseteq}\Hit(\Ccal)
\]
is the blocker, then the inclusion-minimal exact localizations satisfy
\begin{equation}
\Mcal_\Phi(S)=b(\Ccal_\Phi(S)).
\label{eq:app-minimal-localizations-blocker-shadow}
\end{equation}
This is classical transversal / decision-reduct structure. It makes explicit why an inclusion-minimal circuit need not be the unique minimal circuit.

\subsection{Uniform empirical defect and the classical \texorpdfstring{$g_3$}{g3} error}

Suppose $S=\{x_1,\ldots,x_n\}$ carries the uniform empirical law and use $0$--$1$ loss. For each receiver value $z$, let
\[
n_{z,y}
:=
\bigl|\{x\in S:Q_K(x)=z,\ \Phi(x)=y\}\bigr|.
\]
The finite Bayes formula gives
\begin{equation}
\delta_{\mathrm{unif}(S)}(K)
=
1-\frac1n\sum_z\max_y n_{z,y}.
\label{eq:app-empirical-defect-g3}
\end{equation}
To make the functional dependency $K\to d$ exact by deleting rows, one may retain at most one decision value inside each $Q_K$-equivalence class. The largest exact subtable therefore keeps $\sum_z\max_y n_{z,y}$ rows. Hence \eqref{eq:app-empirical-defect-g3} is exactly the classical $g_3$ deletion error for an approximate functional dependency \citep{kivinen1995approximate}.

Variable-precision and decision-theoretic rough-set models likewise study probabilistic or loss-sensitive relaxations of exact reducts \citep{ziarko1993variable,zhao2008attribute}. The paper's distributional content lies upstream of this finite empirical shadow: the probability state is transported through one neural process and receiver refinement is compared with context invariance before these coarse-grainings are taken.
